%% LINELET: líneas de nodos normales a la pared. DEFLATED_CG: grupos de nodos para el problema grueso.
\begin{tikzpicture}
  % capas prismáticas junto a la pared
  \begin{scope}
    \fill[tinta!20] (0,-0.12) rectangle (3.4,0); \node[font=\tiny,text=tenue,anchor=north] at (1.7,-0.12) {pared};
    \foreach \y in {0,0.06,0.14,0.25,0.4,0.6,0.86,1.2} \draw[tenue!70,line width=0.4pt] (0,\y) -- (3.4,\y);
    \foreach \x in {0,0.425,...,3.4} \draw[tenue!70,line width=0.4pt] (\x,0) -- (\x,1.2);
    \draw[ingTau,line width=1.6pt] (1.7,0) -- (1.7,1.2);
    \foreach \y in {0,0.06,0.14,0.25,0.4,0.6,0.86,1.2} \fill[ingTau] (1.7,\y) circle (1.3pt);
    \node[font=\scriptsize,text=ingTau,anchor=west] at (1.8,1.38) {un linelet};
  \end{scope}
  % grupos de la deflación
  \begin{scope}[shift={(4.4,0)}]
    \fill[ingEsq!25] (0,0) -- (1.3,0) -- (1.1,0.7) -- (0,0.8) -- cycle;
    \fill[ingTau!25] (1.3,0) -- (2.6,0) -- (2.7,0.6) -- (1.1,0.7) -- cycle;
    \fill[ingSub!30] (0,0.8) -- (1.1,0.7) -- (1.5,1.4) -- (0,1.4) -- cycle;
    \fill[ingLes!25] (1.1,0.7) -- (2.7,0.6) -- (2.7,1.4) -- (1.5,1.4) -- cycle;
    \draw[tinta,line width=0.5pt] (0,0) rectangle (2.7,1.4);
    \foreach \p in {(0.55,0.4),(1.95,0.33),(0.65,1.1),(2.1,1.0)} \fill[tinta] \p circle (1.6pt);
    \node[font=\scriptsize,anchor=north,align=center] at (1.35,-0.05) {grupos $\to$ problema grueso};
  \end{scope}
\end{tikzpicture}
